\documentclass{article}

\usepackage{tabularx}
\usepackage{booktabs}
\usepackage{longtable}

\title{Reflection and Traceability Report on \progname}

\author{\authname}

\date{}

\input{../Comments} 

\input{../Common}

\begin{document}

\maketitle

\section{Changes in Response to Feedback}

\subsection{SRS and Hazard Analysis}
Table \ref{tab:srs-feedback-tracking} details the feedback received on the SRS 
document, along with the changes implemented to address them and the issues used to track the changes made.
The feedback focused on ensuring every section in the SRS is written with a high degree 
of detail and meeting the course expectations. For example, ensuring the source of constraints 
in the system is clearly communicated through the documentation. 

Furthermore, based on feedback from peer reviews and the teaching assistant, the Hazard Analysis
document was refined to improve clarity and completeness. The wording throughout the
document was revised to ensure that hazards were described precisely and consistently.
This improved the readability of the document and reduced ambiguity in hazard
descriptions.

Feedback also indicated that the backend server component required clearer definition.
This was addressed by expanding the system components section to explicitly describe
the backend server responsibilities, including input validation, communication between
modules, and error handling. This clarification improved the understanding of how failures
within the backend could impact overall system safety.

Additionally, feedback highlighted missing assumptions related to network stability.
A network stability assumption was added to the Critical Assumptions section to reflect
the system’s dependence on reliable communication between components. Including this
assumption helped capture hazards related to connectivity issues and ensured that system
limitations were properly documented.

These revisions strengthened the hazard identification process and improved traceability
between assumptions, system components, and potential risks. The details of the changes 
are outlined in Table \ref{tab:HA-feedback-tracking}.


\renewcommand{\arraystretch}{1.5}
\begin{longtable}{|p{2cm}|p{3.5cm}|p{4.5cm}|p{3cm}|}
    \caption{SRS Feedback Tracking and Implementation} \label{tab:srs-feedback-tracking} \\
    \hline
    \textbf{Feedback Source} & \textbf{Feedback Summary} & \textbf{Changes Made} & \textbf{Related Commits/Issues} \\
    \hline
    \endfirsthead
    
    \multicolumn{4}{c}{{\bfseries Table \thetable\ continued from previous page}} \\
    \hline
    \textbf{Feedback Source} & \textbf{Feedback Summary} & \textbf{Changes Made} & \textbf{Related Commits/Issues} \\
    \hline
    \endhead
    
    \hline \multicolumn{4}{|r|}{{Continued on next page}} \\ \hline
    \endfoot
    
    \hline

    Peer Review & Missing additional BUCs &
Added partition-specific business use cases. &
\href{https://github.com/4G06-CAPSTONE-2025/Reading4All/issues/113}{Issue \#113} \\ \hline

Peer Review & Single BUC too broad &
Decomposed into multiple supporting BUCs. &
\href{https://github.com/4G06-CAPSTONE-2025/Reading4All/issues/118}{Issue \#118} \\ \hline

    Peer Review & Solution introduced in problem statement &
Refined problem statement to focus strictly on the problem. &
\href{https://github.com/4G06-CAPSTONE-2025/Reading4All/issues/26}{Issue \#26} \\ \hline

Peer Review & Goals not clearly measurable &
Added quantitative metrics to make goals measurable. &
\href{https://github.com/4G06-CAPSTONE-2025/Reading4All/issues/27}{Issue \#27} \\ \hline
    \endlastfoot
    
TA & Source of constraints was unclear  & Clarified the source of the constraints. & \href{https://github.com/4G06-CAPSTONE-2025/Reading4All/issues/644}{Issue \#644} \\ \hline
TA & Business rules section was misunderstood & rewrote section as system constraints and domain rules & \href{https://github.com/4G06-CAPSTONE-2025/Reading4All/issues/638}{Issue \#638} \\ \hline
TA & Include more than course deliverables in schedule constraints & goes beyond course deadlines, shows when requirements are implemented & \href{https://github.com/4G06-CAPSTONE-2025/Reading4All/issues/646}{Issue \#646} \\ \hline
TA & Tables not labeled/referenced and notation inconsistencies & Added explicit table references, corrected numbering, and ensured consistent notation across document. & \href{https://github.com/4G06-CAPSTONE-2025/Reading4All/issues/641}{Issue \#641} \\ \hline
TA & Formatting and style issues (spelling, table borders, missing section intros, bitmap figures) & Converted tables to longtable/booktabs style, fixed Canadian spelling and capitalization, added introductory text to Sections 4 and 5, and replaced bitmap figures with PDFs. & \href{https://github.com/4G06-CAPSTONE-2025/Reading4All/issues/640}{Issue \#640} \\ \hline
TA & Missing table for traceability & Added in matrix for traceability & \href{https://github.com/4G06-CAPSTONE-2025/Reading4All/issues/643}{Issue \#643}
\end{longtable}


\renewcommand{\arraystretch}{1.5}
\begin{longtable}{|p{2cm}|p{3.5cm}|p{4.5cm}|p{3cm}|}
    \caption{HA Feedback Tracking and Implementation} \label{tab:HA-feedback-tracking} \\
    \hline
    \textbf{Feedback Source} & \textbf{Feedback Summary} & \textbf{Changes Made} & \textbf{Related Commits/Issues} \\
    \hline
    \endfirsthead
    
    \multicolumn{4}{c}{{\bfseries Table \thetable\ continued from previous page}} \\
    \hline
    \textbf{Feedback Source} & \textbf{Feedback Summary} & \textbf{Changes Made} & \textbf{Related Commits/Issues} \\
    \hline
    \endhead
    
    \hline \multicolumn{4}{|r|}{{Continued on next page}} \\ \hline
    \endfoot
    
    \hline

    Peer Review & Mitigation does not address what happens if the risk occurs &
Added fallback strategies to clarify what actions will be taken if a risk materializes. &
\href{https://github.com/4G06-CAPSTONE-2025/Reading4All/issues/25}{Issue \#25} \\ \hline

TA Feedback & Mitigation for offensive or biased text was too vague &
Specified additional filtering and checking before generated text is shown to the user. &
\href{https://github.com/4G06-CAPSTONE-2025/Reading4All/issues/152}{Issue \#152} \\ \hline

TA Feedback & Data deletion mitigation needed more detail &
Clarified deletion-failure handling by adding logging and administrator notification. &
\href{https://github.com/4G06-CAPSTONE-2025/Reading4All/issues/153}{Issue \#153} \\ \hline

TA Feedback & Missing hazard related to model degradation over time &
Added a hazard and mitigation addressing performance decay as new diagram types emerge. &
\href{https://github.com/4G06-CAPSTONE-2025/Reading4All/issues/154}{Issue \#154} \\ \hline
    \endlastfoot
    
TA & Grammar should be more precise & Reviewed grammar and improvements. & \href{https://github.com/4G06-CAPSTONE-2025/Reading4All/issues/649}{Issue \#649} \\
\hline
Peer Review & Backend Server needs to be more precise & clarify backend server in section 3.3 & \href{https://github.com/4G06-CAPSTONE-2025/Reading4All/issues/155}{Issue \#155} \\
\hline
Peer Review & Missing network stability assumption & revise assumptions in section 4 of HA & \href{https://github.com/4G06-CAPSTONE-2025/Reading4All/issues/151}{Issue \#151} \\
\hline
Peer Review & Missing network stability assumption & revise assumptions in section 4 of HA & \href{https://github.com/4G06-CAPSTONE-2025/Reading4All/issues/151}{Issue \#151} \\
\hline
\end{longtable}

\subsection{Module Guide and Module Interface Specification}


\renewcommand{\arraystretch}{1.5}
\begin{longtable}{|p{2cm}|p{3.5cm}|p{4.5cm}|p{3cm}|}
    \caption{MIS Feedback Tracking and Implementation} \label{tab:DD-feedback-tracking} \\
    \hline
    \textbf{Feedback Source} & \textbf{Feedback Summary} & \textbf{Changes Made} & \textbf{Related Commits/Issues} \\
    \hline
    \endfirsthead
    
    \multicolumn{4}{c}{{\bfseries Table \thetable\ continued from previous page}} \\
    \hline
    \textbf{Feedback Source} & \textbf{Feedback Summary} & \textbf{Changes Made} & \textbf{Related Commits/Issues} \\
    \hline
    \endhead
    
    \hline \multicolumn{4}{|r|}{{Continued on next page}} \\ \hline
    \endfoot
    
    \hline
    \endlastfoot
    
Peer Review & Not all modules in MG match the MIS & added missing module definitions & \href{https://github.com/4G06-CAPSTONE-2025/Reading4All/issues/241}{Issue \#241} \\
\hline
Peer Review & Misalignment of \texttt{userinteface-} \texttt{interactions} module in MIS and MG & clarify module in MIS to reflect MG  & \href{https://github.com/4G06-CAPSTONE-2025/Reading4All/issues/239}{Issue \#239} \\
\hline
TA & Mathematics appendix belonged in the Detailed Design body rather than an appendix & Moved the mathematics content into the appropriate section of the Detailed Design document. & \href{https://github.com/4G06-CAPSTONE-2025/Reading4All/issues/652}{Issue \#652} \\ \hline
TA & Syntax issues in Section 6: optional fields undefined, incorrect output/exception notation, and capitalization of variable names & Defined initialization values for optional fields, updated notation to use \texttt{out :=} and \texttt{exc :=} conventions, and corrected variable name capitalization throughout. & \href{https://github.com/4G06-CAPSTONE-2025/Reading4All/issues/657}{Issue \#657} \\ \hline
TA & Multiple syntax issues across modules: missing environment variables, incorrect types, absent constructors, and misnamed or missing modules & Added environment variables for input/output image files and model storage, corrected variable naming conventions, added constructor and initialization methods to abstract objects, clarified sequence types, used appropriate data types (e.g.\ UUID for user IDs), and added the missing \texttt{FeedbackLoop} module section. & \href{https://github.com/4G06-CAPSTONE-2025/Reading4All/issues/655}{Issue \#655} \\ \hline
TA & CPU and GPU environment variables declared but never referenced in semantics & Updated semantics in affected modules to reference the \texttt{cpu} and \texttt{gpu} environment variables where appropriate. & \href{https://github.com/4G06-CAPSTONE-2025/Reading4All/issues/659}{Issue \#659} \\ \hline
TA & Semantics lacked formality: informal descriptions, missing exception conditions, and improper notation for transitions and conditionals & Decomposed \texttt{filterInput} into local functions, replaced informal descriptions with formal notation (e.g.\ \texttt{clipboard := allText}), and specified exception-raising conditions using proper conditional notation. & \href{https://github.com/4G06-CAPSTONE-2025/Reading4All/issues/661}{Issue \#661} \\ \hline
\end{longtable}


\renewcommand{\arraystretch}{1.5}
\begin{longtable}{|p{2cm}|p{3.5cm}|p{4.5cm}|p{3cm}|}
    \caption{MG Feedback Tracking and Implementation} \label{tab:MG-feedback-tracking} \\
    \hline
    \textbf{Feedback Source} & \textbf{Feedback Summary} & \textbf{Changes Made} & \textbf{Related Commits/Issues} \\
    \hline
    \endfirsthead
    
    \multicolumn{4}{c}{{\bfseries Table \thetable\ continued from previous page}} \\
    \hline
    \textbf{Feedback Source} & \textbf{Feedback Summary} & \textbf{Changes Made} & \textbf{Related Commits/Issues} \\
    \hline
    \endhead
    
    \hline \multicolumn{4}{|r|}{{Continued on next page}} \\ \hline
    \endfoot
    
    \hline
    \endlastfoot
    
Peer Review & Missing Use Hierarchy Diagram & added a use hierarchy diagram to document& \href{https://github.com/4G06-CAPSTONE-2025/Reading4All/issues/234}{Issue \#234} \\
TA & Quality information and formatting inconsistencies in SoftArch section & Standardized AC wording, fixed spacing, corrected Canadian spelling, clarified module types, and ensured hierarchy diagram usage. & \href{https://github.com/4G06-CAPSTONE-2025/Reading4All/issues/650}{Issue \#650} \\ \hline

TA & Presentation issues including table overflow, figure format, and missing references & Converted tables to booktabs style, fixed overflow, added table references in text, and replaced bitmap figures with vector PDF images. & \href{https://github.com/4G06-CAPSTONE-2025/Reading4All/issues/653}{Issue \#653} \\ \hline
\end{longtable}

\subsection{VnV Plan and Report}

Table \ref{tab:VnVPlan-feedback-tracking} details the feedback received from both the teaching assistant and peer reviewers on the 
VnV Plan document which led to several important clarifications 
and improvements to the document. A key change was the addition of clearly defined evaluation metrics for system tests, ensuring 
that requirements such as FR-ST2 are properly aligned with the evaluation framework. Structural improvements were 
also made, including the addition of next steps in the NFR section and the clarification of missing definitions, such as formally 
defining a session and establishing a maximum number of history entries. Additionally, a plan for tracking GitHub issues was outlined 
to improve traceability and organization throughout development. Some feedback items, such as additional checklists and previously 
defined functional requirements, were addressed through issue comments or identified as already covered in existing sections. Overall, 
these updates improved the clarity, completeness, and traceability of the VnV Plan, ensuring it more effectively supports validation 
and verification activities.

\renewcommand{\arraystretch}{1.5}
\begin{longtable}{|p{2cm}|p{3.5cm}|p{4.5cm}|p{3cm}|}
    \caption{VnV Plan Feedback Tracking and Implementation} \label{tab:VnVPlan-feedback-tracking} \\
    \hline
    \textbf{Feedback Source} & \textbf{Feedback Summary} & \textbf{Changes Made} & \textbf{Related Commits/Issues} \\
    \hline
    \endfirsthead
    
    \multicolumn{4}{c}{{\bfseries Table \thetable\ continued from previous page}} \\
    \hline
    \textbf{Feedback Source} & \textbf{Feedback Summary} & \textbf{Changes Made} & \textbf{Related Commits/Issues} \\
    \hline
    \endhead
    
    \hline \multicolumn{4}{|r|}{{Continued on next page}} \\ \hline
    \endfoot
    
    \hline
    \endlastfoot
    
Peer Review & Specify the metric used in FR-ST2 system test  & specified evaluation metrix in appendix 6.1 for FR-ST2 & \href{https://github.com/4G06-CAPSTONE-2025/Reading4All/issues/149}{Issue \#149} \\
\hline
Peer Review & Addition of additional checklists & none; addressed in issue comments & \href{https://github.com/4G06-CAPSTONE-2025/Reading4All/issues/148}{Issue \#148} \\
\hline
Peer Review & Addition of next steps in the NFR section & added next steps in NFR section & \href{https://github.com/4G06-CAPSTONE-2025/Reading4All/issues/146}{Issue \#146} \\
\hline
Peer Review & Missing a definition of a session and maximum number of history items &  A session was defined in document and history limit of 10 entries was set & \href{https://github.com/4G06-CAPSTONE-2025/Reading4All/issues/150}{Issue \#150} \\
\hline
Peer Review & Missing a plan for tracking created GitHub Issues & Outlined a plan for how GitHub Issues will be used & \href{https://github.com/4G06-CAPSTONE-2025/Reading4All/issues/147}{Issue \#147} \\
\hline
Peer Review & Possible missing functional requirement for viewing history & addressed in issue as this is already described in FR 5 & \href{https://github.com/4G06-CAPSTONE-2025/Reading4All/issues/114}{Issue \#114} \\
TA & Traceability coverage unclear & Added traceability matrix and explicit coverage description linking tests to all FRs and NFRs. & \href{https://github.com/4G06-CAPSTONE-2025/Reading4All/issues/647}{Issue \#647} \\ \hline

TA & Feasibility extras not clearly explained & Added description of AI/ML Report and User Testing Report extras and explained their impact on project feasibility. & \href{https://github.com/4G06-CAPSTONE-2025/Reading4All/issues/648}{Issue \#648} \\ \hline

TA & Non-functional requirement tests too ambiguous & Refined NFR tests with clearer acceptance criteria, thresholds, and references to accessibility standards. & \href{https://github.com/4G06-CAPSTONE-2025/Reading4All/issues/645}{Issue \#645} \\ \hline

TA & Spelling, grammar, and formatting inconsistencies & Added introductory sentences, fixed Canadian spelling, improved spacing, and converted tables to booktabs style. & \href{https://github.com/4G06-CAPSTONE-2025/Reading4All/issues/639}{Issue \#639} \\ \hline

Peer Review & Traceability tables would benefit from requirement columns in test result tables and matrix-style reformatting & Added requirement columns to test result tables and reformatted traceability tables as matrices to improve coverage visibility at a glance. & \href{https://github.com/4G06-CAPSTONE-2025/Reading4All/issues/601}{Issue \#601} \\ \hline

Peer Review & Traceability section incomplete: NFR-ST25 and NFR-ST26 missing from both traceability tables & Updated both the requirements trace and module trace tables to include NFR-ST25 and NFR-ST26. & \href{https://github.com/4G06-CAPSTONE-2025/Reading4All/issues/596}{Issue \#596} \\ \hline

Peer Review & Failed tests (FR-ST7, NFR-ST14, NFR-ST25) lacked action plans or hotfix strategies & Documented resolution of FR-ST7 session timeout fix and NFR-ST25 deployment issue under Changes Due to Testing; noted NFR-ST14 documentation availability plan. & \href{https://github.com/4G06-CAPSTONE-2025/Reading4All/issues/594}{Issue \#594} \\ \hline

Peer Review & Failed tests lacked root cause analysis and Section 8 did not map to specific failed test IDs & Added analysis sections beneath failed test tables identifying root causes and linking Section 8 changes to specific failed test IDs. & \href{https://github.com/4G06-CAPSTONE-2025/Reading4All/issues/592}{Issue \#592} \\ \hline

TA & VnV Plan did not specify target scores for survey questions & Added expected target scores for each survey question in the data collection section. & \href{https://github.com/4G06-CAPSTONE-2025/Reading4All/issues/651}{Issue \#651} \\ \hline

TA & Caption scores reported with no improvement plan and no examples of good/bad captions & Added a plan to improve low-scoring captions and included examples of generated captions with their corresponding scores. & \href{https://github.com/4G06-CAPSTONE-2025/Reading4All/issues/656}{Issue \#656} \\ \hline

TA & NFR pass/fail claims lacked supporting data, particularly for accessibility testing & Added quantitative data to support NFR results, including accessibility tool output from Section 9. & \href{https://github.com/4G06-CAPSTONE-2025/Reading4All/issues/658}{Issue \#658} \\ \hline

TA & Changes in response to VnV did not clearly explain how lacking caption areas would be improved & Clarified the improvement plan for caption areas with lower scores in the Changes Due to Testing section. & \href{https://github.com/4G06-CAPSTONE-2025/Reading4All/issues/660}{Issue \#660}
\end{longtable}

Table \ref{tab:vnvReport-feedback-tracking} summarizes the feedback received from peer reviewers and the teaching assistant on the VnV Report, along with the corresponding improvements made to the document. Several updates focused on improving clarity and traceability, including adding introductions to sections, referencing all tables and figures, and correcting formatting inconsistencies. Performance-related feedback led to the definition of symbolic constants and clearer specification of thresholds for alt-text generation and UI responsiveness. Additional analysis was incorporated for previously failed tests, and mitigation plans were documented, including fixes for session timeout functionality and deployment-related issues. Structural enhancements such as improved code coverage explanations, clarification of mutation testing decisions, and corrections to symbolic constants further strengthened the report. Overall, these changes improved the completeness, readability, and traceability of the VnV Report, ensuring that feedback was clearly addressed and documented.
\renewcommand{\arraystretch}{1.5}
\begin{longtable}{|p{2cm}|p{3.5cm}|p{4.5cm}|p{3cm}|}
    \caption{VnV Report Feedback Tracking and Implementation} \label{tab:vnvReport-feedback-tracking} \\
    \hline
    \textbf{Feedback Source} & \textbf{Feedback Summary} & \textbf{Changes Made} & \textbf{Related Commits/Issues} \\
    \hline
    \endfirsthead
    
    \multicolumn{4}{c}{{\bfseries Table \thetable\ continued from previous page}} \\
    \hline
    \textbf{Feedback Source} & \textbf{Feedback Summary} & \textbf{Changes Made} & \textbf{Related Commits/Issues} \\
    \hline
    \endhead
    
    \hline \multicolumn{4}{|r|}{{Continued on next page}} \\ \hline
    \endfoot
    
    \hline
    \endlastfoot
    
Peer Review & NFR section is missing an introduction to outline purpose, scope, etc.   & Added a more detailed introduction to the section & \href{https://github.com/4G06-CAPSTONE-2025/Reading4All/issues/604}{Issue \#604} \\
\hline
TA & No plan in place for if roles need to be switched. & Outlined a plan for how roles were switched when neeeded  & \href{https://github.com/4G06-CAPSTONE-2025/Reading4All/issues/642}{Issue \#642} \\
\hline
Peer Review & Code Coverage needs to be improved and explanation needed for missing coverage. & Increased coverage in files with lower coverage and added an explanation &\href{https://github.com/4G06-CAPSTONE-2025/Reading4All/issues/603}{Issue \#603} \\
\hline 
Peer Review & Failed tests were reported but not analyzed and lacked explanation of why modules/tests failed   & added analysis sections under the tables of where tests failed  & \href{https://github.com/4G06-CAPSTONE-2025/Reading4All/issues/597}{Issue \#597} \\
\hline
Peer Review & Mutation testing was mentioned in the VnV Plan but not discussed in the VnV Report   & Added a subsection explaining why mutation testing was not implemented & \href{https://github.com/4G06-CAPSTONE-2025/Reading4All/issues/600}{Issue \#600} \\
\hline
Peer Review & Symbolic constants not defined and typos  & Added the symbolic constants table and fixed typos & \href{https://github.com/4G06-CAPSTONE-2025/Reading4All/issues/605}{Issue \#605} \\
\hline

Peer Review & Tables/figures not referenced, captions placement, wide tables, and missing section introductions & 
Referenced all tables in text, moved captions above tables, resized tables to avoid margin overflow, and added introductory sentences to sections/subsections &
\href{https://github.com/4G06-CAPSTONE-2025/Reading4All/issues/654}{Issue \#654} \\
\hline

Peer Review & Unclear performance thresholds and contrast testing specification for NFR-ST1 & 
Added symbolic constants defining performance thresholds (T\_ALT\_GEN\_SMALL, T\_ALT\_GEN\_LARGE, T\_UI\_RESP) and clarified NFR-ST1 testing environment limitations for quantitative measurement &
\href{https://github.com/4G06-CAPSTONE-2025/Reading4All/issues/602}{Issue \#602} \\
\hline

Peer Review & Missing mitigation plans for failed tests including FR-ST7 and NFR-ST25 & 
Session timeout failure (FR-ST7) was fixed during testing and documented under Changes Due to Testing; deployment issues (NFR-ST25) were resolved prior to final demo &
\href{https://github.com/4G06-CAPSTONE-2025/Reading4All/issues/594}{Issue \#594} \\
\hline
\end{longtable}

\section{Challenges and Extras}

The challenge level of the project is set at a \textit{general} level, as agreed on with the professor. The project involved integrating multiple technical components, including developing an AI-based system for generating alternative text for complex physics diagrams and designing an accessible user interface. Although these tasks required careful design and coordination across different system components, the overall scope and complexity remained consistent with a general challenge level rather than advanced.

The extras completed for this project include a Machine Learning (ML) Report and a User Testing Report. Both reports provide 
additional insight into the development and evaluation of the Reading4All system. The ML Report documents the design, 
experimentation, and implementation of the machine learning component, including model selection, dataset 
preparation, and iterative improvements across revisions.

The User Testing Report focuses on evaluating the usability and effectiveness of the system from a user 
perspective. It outlines the goals of user testing, the rationale behind the tasks assigned to 
participants, and the methodology used to assess performance. The report also includes supervisor 
evaluation, as well as a detailed analysis of the results obtained from testing. These findings were 
used to identify usability issues and guide improvements to both the system and the generated alt text.

\section{Design Iteration (LO11 (PrototypeIterate))} 

Our design has evolved from the first version in Rev0 to Rev1. Firstly, looking at our user interface changes, 
many changes were made to improve the usability and accessibility of the \texttt{Reading4All} application. \\

The first improvement came from our VNV findings, which showed that the \texttt{View History} button was not clearly recognizable as button.
We updated its styling to use a maroon color and stronger visual cues so its more easily identifiable as a clickable button.
Secondly, based on feedback from our supervisor, we made the maximum file size clearly noted to users before uploading. 
This helps users understand the limit upfront and reduces potential errors.
The third improvement also came from supervisor feedback, which was to make the ordering of images in the history page more clear to users. As such, we stated to 
users in the history page that entires are shown \texttt{with the most recent shown first}. 
The final UI improvement came from our VnV findings, which highlighted issues with text readability after alt text generation, 
specifically words being split across lines. To address this, we updated the text wrapping  to ensure words are no longer broken between lines and remain easy to read.
Ultimately, these design changes were guided from supervisor feedback and testing, highlighting issues that were not initially obvious to the development team.\\

Furthermore, several changes were made to improve the quality of the alt text generated by the \texttt{Reading4All} application. 
Firstly, based on feedback from the course instructor regarding the alt text quality and overall flow, we implemented a post-processing layer. 
After the BLIP AI model generates the alt text, it is passed through a function that fixes formatting issues to improve readability and overall presentation.
Secondly, we addressed feedback from our TA and supervisor regarding the quality of the alt text generated during Rev0. 
To improve this, we expanded our training dataset by adding 200 additional annotated images. 
This helped the model generate more accurate descriptions and overall quality.\\

Overall, the design evolved from an initially prototype to a more refined and intuitive system, with improvements 
focused on usability, accessibility and alt text quality based on testing, supervisor, instructor and TA feedback. 

\section{Design Decisions (LO12)}

Several frontend design decisions were made to improve accessibility and usability. 
We designed the interface using appropriate contrast ratios and colour selections 
based on WCAG and AODA accessibility guidelines. This decision was made to ensure 
that the application is usable by individuals with visual impairments and complies 
with accessibility standards. Following these guidelines also improves readability 
and overall user experience.

Another key decision was implementing a tab-based navigation system for the frontend. 
This structure allows users to easily switch between different functionalities without 
overwhelming the interface. The tab-based layout improves organization and reduces 
cognitive load, making the system more intuitive to navigate.

We also decided to include a session history feature to allow users to view previous 
interactions. This functionality helps users track past generated content and improves 
workflow efficiency by allowing them to reference earlier results without repeating 
actions.

Additional functionalities such as editing alt text and copying generated output were 
added later in the design process. These features were introduced based on usability 
considerations to give users more control over the generated content and make it easier 
to reuse results in other applications. These improvements enhanced flexibility and 
overall usability of the system.

\section{Economic Considerations (LO23)}

The Reading4All tool is designed to improve accessibility for physics diagrams by automatically generating alternative text. 
The system is developed with accessibility standards in mind, including WCAG~2.1 and AODA compliance, and supports 
tab-based keyboard navigation. Creating high-quality alternative text for complex STEM diagrams is often difficult and 
time-consuming, particularly when balancing appropriate length, level of detail, and contextual accuracy. This tool helps 
students, teaching assistants, instructors, and academic departments quickly generate descriptive alternative text, providing 
a strong starting point that can be refined as needed.

Currently, the tool is intended for use within McMaster University. The primary users include students with accessibility needs, 
instructors preparing course materials, and teaching assistants responsible for ensuring content accessibility. Although the 
tool is not intended for commercial sale, there is clear institutional value in improving accessibility compliance and reducing 
manual effort required to generate alternative text.

The system relies on cloud-based infrastructure including Vercel for frontend deployment, Render for backend services, Supabase 
for database and authentication, and Hugging Face for hosting machine learning models. The estimated operational costs depend 
on usage scale. For a small institutional deployment, hosting costs are expected to remain relatively low, with frontend hosting 
typically ranging from \$0--\$20 per month, backend hosting from \$10--\$40 per month, database and storage costs from 
\$0--\$25 per month, and model hosting or inference costs ranging from \$0--\$50 per month depending on usage.

The primary development-related costs would involve expanding datasets and improving model quality. Collecting and curating 
physics diagram datasets, including sourcing open-source images, may cost approximately \$200--\$800. Subject matter experts 
to verify and annotate diagrams could require an estimated \$300--\$1000 depending on the number of images reviewed. Additional 
costs for user surveys, accessibility testing, and evaluation efforts may range from \$100--\$300. Optional performance upgrades, 
such as faster inference or GPU-based model hosting, could increase operational costs by approximately \$20--\$100 per month.

Because this project is intended to remain open source and institutionally supported, revenue generation is not the primary goal. 
Instead, adoption would occur through accessibility initiatives, course integration, and departmental recommendations. The 
potential user base includes students requiring accessible materials, instructors creating lecture content, and accessibility 
offices across McMaster University. Future expansion to other academic institutions could further increase the number of potential 
users while maintaining relatively low operational costs.

\subsection{How Does Your Project Management Compare to Your Development Plan}

Overall, the team mostly followed the Development Plan throughout the project, with 
some deviations that emerged as the project progressed. The team meeting plan was 
largely adhered to, with regular meetings held consistently throughout the term to 
discuss progress, assign tasks, and address blockers. Communication was primarily 
conducted through Discord, which aligned with the plan and proved effective for 
day-to-day coordination and quick decision-making.

GitHub Issues and Git were used as the primary workflow tools, consistent with what 
was outlined in the Development Plan. Issues were created to track feedback, 
document changes, and assign work to team members. This approach supported 
traceability across documents and implementation tasks. The technologies used during 
development, including React for the frontend, a cloud-based backend, and Hugging 
Face for model hosting, aligned with the planned technology stack.

The main deviations from the plan arose from underestimating the complexity of 
certain tasks, particularly around machine learning model refinement and system 
integration. These complexities required more time and iteration than originally 
anticipated, which occasionally affected the pacing of deliverables relative to the 
planned schedule.

\subsection{What Went Well?}

Team communication was one of the strongest aspects of the project management 
process. The use of Discord allowed the team to stay aligned, share updates quickly, 
and resolve issues without waiting for scheduled meetings. This reduced delays and 
kept development momentum consistent throughout the term.

Regular team meetings also contributed positively to the project. Holding consistent 
meetings ensured that all members were aware of current priorities, upcoming 
deadlines, and any blockers that needed to be addressed. These meetings created a 
shared understanding of progress and helped distribute work effectively across the 
team.

GitHub Issues provided a structured way to track feedback, assign tasks, and 
document decisions. Using issues to link feedback to specific commits and document 
changes improved traceability across all project deliverables and made it easier to 
demonstrate that feedback had been addressed.

\subsection{What Went Wrong?}

The primary challenge encountered during the project was underestimating the 
complexity of several tasks. In particular, training and refining the machine learning 
model required significantly more iteration than planned, including expanding the 
dataset and implementing post-processing steps to improve alt text quality. This 
affected the time available for other tasks and required the team to re-prioritize 
work at certain points.

Integration and merging issues also caused friction during development. As multiple 
team members worked on different components simultaneously, combining changes 
occasionally introduced conflicts that required additional time to resolve. This was 
particularly evident during periods of active development when several features were 
being implemented in parallel.

\subsection{What Would you Do Differently Next Time?}

In future projects, the team would invest more time upfront in breaking down complex 
tasks into smaller, well-scoped subtasks before beginning implementation. This would 
produce more realistic time estimates and reduce the risk of underestimating effort 
for technically challenging components such as machine learning pipelines.

To address integration friction, the team would adopt a more disciplined branching 
strategy from the outset, with clearer guidelines around when to merge and how to 
handle conflicts. Establishing short integration cycles and more frequent merges 
would help catch conflicts earlier and reduce the overhead of resolving large 
divergences late in the development process.

\section{Ethics and Equity (LO18)}

The \texttt{Reading4All} team applied equity and universal design principles 
throughout the development process to ensure the system is accessible and fair to 
all stakeholders.

From a user interface perspective, the application was designed in compliance with 
WCAG~2.1 and AODA accessibility guidelines. Appropriate colour contrast ratios were 
maintained throughout the interface, and tab-based keyboard navigation was 
implemented to ensure the application is usable by individuals who rely on keyboard 
access or assistive technologies. These decisions were made deliberately to ensure 
that the tool itself does not create additional barriers for the users it is intended 
to support.

The system was designed with users who have visual impairments as a primary 
stakeholder group. Rather than treating accessibility as an afterthought, the core 
purpose of the tool, generating alternative text for physics diagrams, directly 
addresses a gap in accessible educational content. By automating the generation of 
descriptive alt text, the system reduces the burden on instructors and accessibility 
offices while improving the quality and availability of accessible materials for 
students who rely on screen readers or other assistive technologies.

Equity considerations also informed the machine learning component of the project. 
During dataset curation, care was taken to include a diverse range of diagram types 
and styles to reduce the risk of the model performing well only on a narrow subset 
of inputs. Expanding the training dataset by 200 additional annotated images during 
Rev1 was partly motivated by the need to improve generalization and reduce bias in 
generated descriptions.

User testing was conducted with participants representing a range of backgrounds and 
familiarity with the system. Feedback from testing was used to identify usability 
issues that were not apparent to the development team, ensuring that design 
improvements reflected the needs of real users rather than assumptions made 
internally. Supervisor evaluation was also incorporated to provide an external 
perspective on system effectiveness and equity of outcomes.

Overall, the team treated equitable access as a foundational design constraint 
rather than an optional feature, ensuring that the system serves its intended 
users effectively and responsibly.

\section{Reflection on Capstone}

\subsection{Which Courses Were Relevant}

Several courses taken throughout the Software Engineering program were directly 
relevant to the work completed during the capstone project.

Software design courses such as SFWRENG 2AA4 provided a foundational framework for 
structuring the system using modular design principles. Concepts such as information 
hiding, module decomposition, and the separation of concerns were applied directly 
in the development of the Module Guide and Module Interface Specification, and 
informed how the system was divided into frontend, backend, and machine learning 
components.

Courses covering machine learning and artificial intelligence were relevant to the 
core technical component of the project. The knowledge gained from these courses 
supported understanding of model selection, training procedures, and evaluation 
metrics, which were applied when fine-tuning the BLIP model and assessing its 
performance on physics diagram datasets.

Database courses informed the design and use of Supabase for user authentication, 
session storage, and history management. Understanding relational data models and 
query design helped the team structure the database schema to support the required 
application features effectively.

Software testing and quality assurance courses were directly applicable to the VnV 
Plan and VnV Report. Concepts such as test case design, traceability matrices, 
code coverage analysis, and non-functional requirement testing shaped the team's 
approach to verification and validation throughout the project.

\subsection{Knowledge and Skills Outside of Courses}

Several skills and areas of knowledge required for the capstone project were not 
covered in depth by the courses taken during the program and had to be acquired 
independently.

Fine-tuning a pre-trained machine learning model, specifically adapting BLIP for 
physics diagram captioning, required learning techniques not covered in coursework, 
including dataset preparation for image-captioning tasks, transfer learning 
workflows, and iterative evaluation of generated output quality. The team developed 
this knowledge through independent study, experimentation, and review of relevant 
documentation and research.

Cloud deployment was another area that required self-directed learning. Deploying 
the frontend on Vercel, the backend on Render, and the machine learning model on 
Hugging Face involved understanding deployment pipelines, environment configuration, 
and managing communication between distributed services. These practical skills were 
not extensively covered in coursework but were essential to delivering a functional, 
hosted system.

Building the frontend using React required familiarity with component-based 
architecture, state management, and responsive design that went beyond what was 
taught in courses. The team developed these skills through hands-on implementation 
and reference to external documentation.

Finally, applying accessibility standards such as WCAG~2.1 and AODA in a practical 
software context required knowledge that extended beyond general course material. 
Understanding specific contrast ratio requirements, keyboard navigation patterns, 
and screen reader compatibility involved independent research and testing to ensure 
the application met the necessary standards.

\end{document}